\documentclass[10pt,a4paper]{amsart}
\usepackage[margin=19mm]{geometry}
\usepackage{amsmath,amssymb,mathtools}
\usepackage[scr=rsfs,cal=euler]{mathalfa}
\usepackage{xfrac}
\usepackage[unicode,hidelinks,bookmarks=false]{hyperref}
\newcommand{\ie}{i.e.\ }
\newcommand{\cf}{cf.\ }
\newcommand{\resp}{respectively\ }
\newcommand{\Subsection}{\subsection}
\providecommand{\nobreakdash}{\nobreak\hbox{-}}
\setlength{\emergencystretch}{2em}
\multlinegap0pt
\usepackage{bm}
\usepackage{bbm}
\usepackage{dsfont}
\usepackage{xspace}
\usepackage{cancel}
\usepackage{mathtools}
\usepackage{amsthm}
\makeatletter
\@ifundefined{lemma}{\newtheorem{lemma}{Lemma}}{}
\@ifundefined{theorem}{\newtheorem{theorem}{Theorem}}{}
\makeatother
\title[Structural and extremal properties of $l\_1$-Fiedler value]{Structural and extremal properties of $l\_1$-Fiedler value}
\author{M. Rajesh Kannan, Rahul Roy}
\date{Source: arXiv:2601.05771v1 (CC BY 4.0)}
\begin{document}
\maketitle

\noindent\emph{Source and licence.} Structural and extremal properties of $l_1$-Fiedler value. Version arXiv:2601.05771v1 (CC BY 4.0), distributed under the Creative Commons Attribution 4.0 International licence, as recorded in the arXiv licence metadata for this exact version and shown on the abstract page https://arxiv.org/abs/2601.05771v1. Full rights notice: licenses/arxiv-2601.05771-CC-BY-4.0.txt.\par\smallskip

\noindent\textit{Editorial scope.} Counted unit: the Theorem whose source label is \texttt{extremal-graph-max-vertex-degree} in the article (source0.tex lines 547--570, source label \texttt{extremal-graph-max-vertex-degree}), delivered together with the article's own complete original proof of it (source0.tex lines 571--763, one \texttt{proof} environment of 193 lines). No mathematics is written, repaired or re-derived: statement and proof are the article's own text line for line. Required context retained uncounted: extended centre edge result (lines 315--321); centre vertex domination (lines 537--539). Every definition, display and label of those lines is retained unchanged. Omitted from this package and not redistributed: 1--314, 322--536, 540--546, 764--1261. Nothing inside a retained excerpt is omitted or altered: every retained body is delivered exactly as the article's lines are written, so no omission receipt of any kind accompanies this package. Citations: the retained excerpts carry no citation command at all, so no bibliography entry of the article is transcribed and no reference is redistributed. Reference adaptation: none was needed and none was made; every reference inside the delivered unit resolves inside it, so no unresolved reference or citation is produced. Printed slips kept literal: no printed word, symbol, hypothesis or number of the counted statement or of its proof was altered. Layout and class: the article's own class is \texttt{article} with packages including geometry, amsfonts, amsmath, amsthm, amssymb, hyperref, fullpage, graphicx, listings, xcolor, enumitem, tikz, bm; the wrapper is the lane's pinned \texttt{amsart} editorial wrapper at 10pt on A4, which restates the theorem-like environments the retained text needs and adds the article's own macro definitions that the retained text needs (none), with extra wrapper packages (bm, bbm, dsfont, xspace, cancel, mathtools). The delivered numbering is the unit's own. Assets: no retained span contains a figure environment, a graphics-inclusion command, a PDF-page inclusion, a table or a TikZ picture; no image file is copied into this package, the package declares no asset dependency and no image file is redistributed with it. Licence: version arXiv:2601.05771v1 of this article is distributed under the Creative Commons Attribution 4.0 International licence, as recorded in the arXiv licence metadata for this exact version and shown on the abstract page https://arxiv.org/abs/2601.05771v1. A full rights notice is delivered at licenses/arxiv-2601.05771-CC-BY-4.0.txt. Characters: every retained excerpt is pure ASCII, so the delivered engine, XeLaTeX with its default Latin Modern text font, reports no missing character.
\begin{theorem}\label{extended centre edge result}
    Let $T=(V, E)$ be a tree. Let $uv \in E$. Then 
    $$
b(T)=\frac{1}{2}\left(\frac{1}{\vert V_u \vert} +\frac{1}{\vert V_v\vert} \right)
    $$
    if and only if $uv$ is a centre edge in $T$.
\end{theorem}

\begin{lemma}\label{centre vertex domination}
    Let $T=(V,E)$ be a tree and $u$ be a star-root vertex of  $T$. If $uv \in E$, then $\vert V_u \vert \geq \vert V_v\vert.$
\end{lemma}

\begin{theorem}\label{extremal-graph-max-vertex-degree}
    Let $T$ be a tree on $n$ vertices with maximum vertex degree $d_{\max}.$       
\begin{enumerate}
    \item If $n$ is even, then
    \[
        b(T) \geq 
        \begin{cases} 
            \frac{2}{n}, & d_{\max} \leq \frac{n}{2}, \\
            \frac{n}{2d_{\max}(n-d_{\max})}, & d_{\max} > \frac{n}{2}.
        \end{cases}
    \]
    \item If $n$ is odd, then
    \[
        b(T) \geq 
        \begin{cases} 
            \frac{2n}{n^2 - 1}, & d_{\max} \leq \lfloor \frac{n}{2} \rfloor + 1, \\
            \frac{n}{2d_{\max}(n-d_{\max})}, & d_{\max} > \lfloor \frac{n}{2} \rfloor + 1.
        \end{cases}
     \]
     \item If $d_{\max}$ divides $n-1$, then $b(T) \leq \frac{n}{2k(n-k)}$ where $n-1=kd_{\max}$ for some $k \in \mathbb{N}$.
     \item If $n-1=k d_{\max}+r$ for some $k,r \in \mathbb{N}$ with $1 \leq r <d_{\max}$, then $b(T) \leq \frac{n}{2(k+1)(n-k-1)}.$
\end{enumerate}
Moreover, all the bounds mentioned above are sharp.
 \end{theorem}

\begin{proof}
Construct a tree $T$ which minimizes $b(T)$ as follows:  Let $v_1$ be a vertex of degree $d_{\max}$, and attach a path on $n-d_{\max}-1$ vertices to one of the neighbors of $v_1$.


\noindent\textbf{Proof of (1): } Let $n=2l$ for some positive integer $l$.
If $d_{\max} \leq n/2$, then there is a sparsest cut induced by exactly ${n/2}$ vertices in $T$. That is, $b(T)=2/n$. So, $T$ minimizes $b(T)$ in this case.
% Here we are giving an example of such a tree with $n=8$ and $d=4$.
% $$    
%     \begin{tikzpicture}[node distance={15mm}, thick, main/.style = {draw, circle}] 
% \node[main] (1) {$v_1$};
% \node[main] (2) [  right of=1] {$v_2$};
% \node[main] (3) [above of=1] {$v_3$}; 
% \node[main] (4) [ left of=1] {$v_4$};
% \node[main] (5) [below of=1] {$v_5$};
% \node[main] (6) [ right of=2] {$v_6$};
% \node[main] (7) [ right of=6] {$v_7$};
% \node[main] (8) [ right of=7] {$v_8$};

% \draw (1) -- (2);
% \draw (1) -- (3);
% \draw (1) -- (4);
% \draw (1) -- (5);
% \draw (2) -- (6);
% \draw (6) -- (7);
% \draw (7) -- (8);

% \end{tikzpicture}
% $$
If $d_{\max} > n/2$, then there is a sparsest cut induced by $n-d_{\max}$ number of vertices in $T$. That is, $b(T)=\frac{n}{2d_{\max}(n-d_{\max})}$.  We claim $T$ minimizes $b(T)$ in this case. Suppose that,  $T^\prime=(V^\prime,E^\prime)$ is a tree with $b(T^\prime) <b(T)$. Let $w^\prime \in V^\prime$ with $\deg(w^\prime)=d_{\max}$. Let $u^\prime v^\prime$ be a centre edge in $T^\prime$. From Theorem \ref{extended centre edge result}, we get
 $$
\frac{n}{2}\frac{1}{\vert V^\prime_{u^\prime}\vert \vert V^\prime_{v^\prime}\vert}=b(T^\prime)<b(T)=\frac{n}{2}\frac{1}{d_{\max}(n-d_{\max})}.
 $$
  Therefore, $\vert V^\prime_{u^\prime}\vert \vert V^\prime_{v^\prime}\vert > d_{\max}(n-d_{\max})$. Without loss of generality, let $\vert V^\prime_{u^\prime}\vert \leq \vert V^\prime_{v^\prime}\vert$. Then $n-d_{\max} < \vert V^\prime_{u^\prime}\vert \leq l <d_{\max}$ and $l \leq \vert V^\prime_{v^\prime}\vert <d_{\max}$. Note that there is exactly one edge between $T^\prime_{u^\prime}$ and $T^\prime_{v^\prime}$, so the subtree that containing $w$ should contain at least $d_{\max}-1$ neighbors of $w$. This implies that at least one of the subtrees should be of order at least $d_{\max}$. That is, either $\vert T^\prime_{u^{'}} \vert \geq d_{\max}$ or  $\vert T^\prime_{v^{'}} \vert \geq d_{\max}$, which is  a contradiction. Thus $T$ minimizes $b(T)$.
% Here we are giving an example of such a tree with $n=8$ and $d=5$. 
% $$    
%     \begin{tikzpicture}[node distance={15mm}, thick, main/.style = {draw, circle}] 
% \node[main] (1) {$v_1$};
% \node[main] (2) [  right of=1] {$v_2$};
% \node[main] (3) [above of=1] {$v_3$}; 
% \node[main] (4) [ left of=1] {$v_4$};
% \node[main] (5) [below of=1] {$v_5$};
% \node[main] (6) [ below right of=1] {$v_6$};
% \node[main] (7) [ right of=2] {$v_7$};
% \node[main] (8) [ right of=7] {$v_8$};

% \draw (1) -- (2);
% \draw (1) -- (3);
% \draw (1) -- (4);
% \draw (1) -- (5);
% \draw (1) -- (6);
% \draw (2) -- (7);
% \draw (7) -- (8);

% \end{tikzpicture}
% $$

\noindent\textbf{Proof of (2): } Let $n=2l+1$ for some positive integer $l$.
 Let $d_{\max} \leq \lfloor \frac{n}{2} \rfloor +1$. Then there is a sparsest cut in $T$ induced by $l$ vertices in $T$ with $b(T)=\frac{2n}{n^2-1}$. Consequently, $T$ minimises $b(T)$. If $d_{\max} > \lfloor \frac{n}{2} \rfloor +1$, then there is a sparsest cut induced by $n-d_{\max}$ vertices in $T$ with $b(T)=\frac{n}{2d_{\max}(n-d_{\max})}$. The remainder of the proof follows along the same lines as in part (1).
%  Here is an example of a tree with $n=9$ and $d=5$.
%  $$    
%     \begin{tikzpicture}[node distance={15mm}, thick, main/.style = {draw, circle}] 
% \node[main] (1) {$v_1$};
% \node[main] (2) [  right of=1] {$v_2$};
% \node[main] (3) [above of=1] {$v_3$}; 
% \node[main] (4) [ left of=1] {$v_4$};
% \node[main] (5) [below of=1] {$v_5$};
% \node[main] (6) [ below right of=1] {$v_6$};
% \node[main] (7) [ right of=2] {$v_7$};
% \node[main] (8) [ right of=7] {$v_8$};
% \node[main] (9) [ right of=8] {$v_9$};

% \draw (1) -- (2);
% \draw (1) -- (3);
% \draw (1) -- (4);
% \draw (1) -- (5);
% \draw (1) -- (6);
% \draw (2) -- (7);
% \draw (7) -- (8);
% \draw (8) -- (9);
% \end{tikzpicture}
% $$

% Here is one example of a tree with $n=9$ and $d=6$.
% $$    
%     \begin{tikzpicture}[node distance={15mm}, thick, main/.style = {draw, circle}] 
% \node[main] (1) {$v_1$};
% \node[main] (2) [  right of=1] {$v_2$};
% \node[main] (3) [above of=1] {$v_3$}; 
% \node[main] (4) [ left of=1] {$v_4$};
% \node[main] (5) [below of=1] {$v_5$};
% \node[main] (6) [ below right of=1] {$v_6$};
% \node[main] (7) [ above left of=1] {$v_7$};
% \node[main] (8) [ right of=2] {$v_8$};
% \node[main] (9) [ right of=8] {$v_9$};

% \draw (1) -- (2);
% \draw (1) -- (3);
% \draw (1) -- (4);
% \draw (1) -- (5);
% \draw (1) -- (6);
% \draw (1) -- (7);
% \draw (2) -- (8);
% \draw (8) -- (9);
% \end{tikzpicture}
% $$


\noindent\textbf{Proof of (3):} Let $d_{\max}$ divides $n-1$, and let $n-1 = k d_{\max}$. Construct the tree $T$ as follows: Let $\deg(v_1)=d_{\max}$. Attach to each pendant neighbor of $v_1$ a branch of size $k-1$, ensuring that the degree of every vertex in the resulting graph is at most $d_{\max}$. This procedure may produce more than one nonisomorphic tree. Fix any one of the trees obtained in this way, and denote it by $T$. 
 
% Here we are giving an example where $n=13$ and $d_{\max}=3$.
% $$    
%     \begin{tikzpicture}[node distance={15mm}, thick, main/.style = {draw, circle}] 
% \node[main] (1) {$v_1$};
% \node[main] (2) [ above right of=1] {$v_2$};
% \node[main] (3) [below of=1] {$v_3$}; 
% \node[main] (4) [ above left of=1] {$v_4$};
% \node[main] (5) [ above right of=2] {$v_5$};
% \node[main] (6) [ below right of=2] {$v_6$};
% \node[main] (7) [ right of=5] {$v_7$};
% \node[main] (8) [ below left of=3] {$v_8$};
% \node[main] (9) [ below right of=3] {$v_9$};
% \node[main] (10) [ below  of=9] {$v_{10}$};
% \node[main] (11) [ above left of=4] {$v_{11}$};
% \node[main] (12) [ below left of=4] {$v_{12}$};
% \node[main] (13) [ left of=12] {$v_{13}$};

% \draw (1) -- (2);
% \draw (1) -- (3);
% \draw (1) -- (4);
% \draw (2) -- (5);
% \draw (2) -- (6);
% \draw (3) -- (8);
% \draw (3) -- (9);
% \draw (4) -- (11);
% \draw (4) -- (12);
% \draw (5) -- (7);
% \draw (9) -- (10);
% \draw (12) -- (13);


% \end{tikzpicture}
% $$
In $T$, the $k$ vertices of any single branch induces a sparsest cut. Thus, $b(T)=\frac{n}{2k(n-k)}$. We now show that $T$ maximizes $b(T)$. Suppose, to the contrary, that there exists a tree  $T^\prime=(V^\prime,E^\prime)$ such that $b(T^\prime)>b(T)$. Let $w^\prime \in V^\prime$ be a star-root vertex of $T^\prime$, and let $w^\prime v^\prime$ be a centre edge in $T^\prime$. By Theorem \ref{extended centre edge result}
$$
 \frac{n}{2}\frac{1}{\vert V^\prime_{w^\prime}\vert \vert V^\prime_{v^\prime}\vert}=b(T^\prime)>b(T)=\frac{n}{2k(n-k)}.
 $$
Hence $\vert V^\prime_{w^\prime}\vert \vert V^\prime_{v^\prime}\vert < k(n-k)$. By Lemma \ref{centre vertex domination}, we have $\vert V^\prime_{w^\prime}\vert \geq \vert V^\prime_{v^\prime}\vert$. It follows that $$\vert V^\prime_{v^\prime}\vert \leq k-1~\mbox{~ and~} \vert V^\prime_{w^\prime} \vert \geq n-k+1.$$ Therefore, $$\vert V^\prime_{w^\prime} \vert - \vert V^\prime_{v^\prime}\vert \geq n-2k+2.$$
Now we claim that every branch of $w^\prime$ holds at most $k-1$ vertices. Suppose that, there exists a branch of $w^\prime$ that contains at least $k$ vertices. Let $z^\prime$ be the neighbor of $w^\prime$ in this branch. By Lemma \ref{centre vertex domination}, we have $\vert V^\prime_{w^\prime}\vert \geq \vert V^\prime_{z^\prime} \vert$. Also  $\vert V^\prime_{w^\prime}\vert \leq n-k$ and $\vert V^\prime_{z^\prime}\vert \geq k$. Therefore $$\vert V^\prime_{w^\prime}\vert-\vert V^\prime_{z^\prime}\vert \leq n-2k <n-2k+2 \leq \vert V^\prime_{w^\prime} \vert - \vert V^\prime_{v^\prime}\vert,$$
which contradicts the assumption that $w^\prime v^\prime$ is a centre edge of $T^\prime$.
Since $deg(w^\prime) \leq d_{\max}$, the number of vertices in $T^\prime$ is at most $d_{\max}(k-1)+1 <d_{\max}k+1=n$, which is a contradiction. Therefore $T$ maximizes $b(T)$.

\noindent\textbf{Proof of (4):} Let $n-1=d_{\max}k+r$, where $k$ and $r$ are non-negative integers with $1 \leq r <d_{\max}$.
In this case, we use a construction for $T$ similar to that in the previous case. In addition, the remaining $r$ vertices are distributed among distinct branches of $v_1$, one vertex per branch. As a result, exactly $r$ branches containing $k+1$ vertices each, while the remaining branches contain exactly $k$ vertices each.
 % Here we are giving an example with $n=15$ and $d_{\max}=3$. So $r=2$.
% $$    
%     \begin{tikzpicture}[node distance={15mm}, thick, main/.style = {draw, circle}] 
% \node[main] (1) {$v_1$};
% \node[main] (2) [ above right of=1] {$v_2$};
% \node[main] (3) [below of=1] {$v_3$}; 
% \node[main] (4) [ above left of=1] {$v_4$};
% \node[main] (5) [ above right of=2] {$v_5$};
% \node[main] (6) [ below right of=2] {$v_6$};
% \node[main] (7) [ right of=5] {$v_7$};
% \node[main] (8) [ below left of=3] {$v_8$};
% \node[main] (9) [ below right of=3] {$v_9$};
% \node[main] (10) [ below  of=9] {$v_{10}$};
% \node[main] (11) [ above left of=4] {$v_{11}$};
% \node[main] (12) [ below left of=4] {$v_{12}$};
% \node[main] (13) [ left of=12] {$v_{13}$};
% \node[main] (14) [ right of=6] {$v_{14}$};
% \node[main] (15) [ below of=8] {$v_{15}$};

% \draw (1) -- (2);
% \draw (1) -- (3);
% \draw (1) -- (4);
% \draw (2) -- (5);
% \draw (2) -- (6);
% \draw (3) -- (8);
% \draw (3) -- (9);
% \draw (4) -- (11);
% \draw (4) -- (12);
% \draw (5) -- (7);
% \draw (6) -- (14);
% \draw (8) -- (15);
% \draw (9) -- (10);
% \draw (12) -- (13);



% \end{tikzpicture}
% $$
 The set of vertices in any branch of $T$ containing $k+1$ vertices induces a sparsest cut in $T$.  Hence, $b(T)=\frac{n}{2(k+1)(n-k-1)}$. Suppose, for contradiction, that there exists a tree $T^\prime=(V^\prime,E^\prime)$ such that $b(T^\prime) >\frac{n}{2(k+1)(n-k-1)}$. Let $w^\prime \in V^\prime$ be a star-root vertex of $T^\prime.$ Arguing as in the previous case, we conclude that every branch of $w^\prime$ has at most $k$ vertices. Hence, $T^\prime$ has at most $kd_{\max}+1<kd_{\max}+1+r=n$ vertices, which is a contradiction. Thus, $T$ maximizes $b(T)$.
\end{proof}

\end{document}
